% Generated draft: requires mathematical and source review.
\documentclass[11pt]{article}
\usepackage[margin=25mm]{geometry}
\usepackage{fontspec}
\setmainfont{Latin Modern Roman}
\usepackage{amsmath,amssymb,mathtools,unicode-math}
\setmathfont{Latin Modern Math}
\usepackage{enumitem,needspace,xurl,hyperref}
\setlength{\parindent}{0pt}
\setlength{\parskip}{5pt}
\setlength{\emergencystretch}{3em}
\newcommand{\N}{\mathbb{N}}
\newcommand{\Z}{\mathbb{Z}}
\newcommand{\Q}{\mathbb{Q}}
\newcommand{\R}{\mathbb{R}}
\newcommand{\C}{\mathbb{C}}
\newcommand{\F}{\mathbb{F}}
\newcommand{\A}{\mathbb{A}}
\newcommand{\PP}{\mathbb{P}}
\newcommand{\E}{\mathbb{E}}
\newcommand{\eps}{\varepsilon}
\newcommand{\dd}{\,\mathrm{d}}
\begin{document}
\section*{UnsolvedMath ID: 20001919}
\textbf{Source problem code:} catalogue entry
\subsection{Definitions and mathematical statement}
Let \(g_{\mathrm{rd}}\) be the unit round metric on \(S^2\), and for a smooth Riemannian metric \(g\) let \(\ell_{\min}(g)\) be the least length of a nonconstant closed geodesic. The source calls this \(\operatorname{sys}(S^2,g)\); it cannot mean a noncontractible loop, since \(S^2\) is simply connected. Babenko's question asks whether every area-\(4\pi\) metric sufficiently close to \(g_{\mathrm{rd}}\) has \(\ell_{\min}(g)\leq2\pi\). The source does not state a topology.

Under the \(C^2\), and hence the \(C^\infty\), interpretation, the supplied aid reports an affirmative theorem of Abbondandolo, Bramham, Hryniewicz and Salomão. More precisely, if a smooth metric on \(S^2\) has positive Gaussian curvature and is \(\delta\)-pinched in the sense
\[
 0<\delta K_{\max}\leq K(x)\leq K_{\max},
 \qquad \frac{4+\sqrt7}{8}<\delta\leq1,
\]
then
\[
 \ell_{\min}(g)^2\leq\pi\operatorname{Area}(S^2,g),
\]
with equality exactly for Zoll metrics in this class. Here Zoll means that all unit-speed geodesics are simple closed geodesics with the same prime length. Smooth \(C^2\)-closeness to \(g_{\mathrm{rd}}\) forces such positive pinching. Thus there is a \(C^2\)-neighborhood \(\mathcal U\) of \(g_{\mathrm{rd}}\) such that
\[
 g\in\mathcal U,\quad\operatorname{Area}(S^2,g)=4\pi
 \quad\Longrightarrow\quad\ell_{\min}(g)\leq2\pi.
\]
The topology is material: no \(C^0\) or Gromov--Hausdorff version is inferred from this statement. Equality in the local theorem need not single out only the round metric; its stated equality class is Zoll.

The aid also reports a separate global special case. If the standard antipodal involution \(a(x)=-x\) is an isometry of \(g\), then without closeness or curvature assumptions
\[
 a^*g=g\quad\Longrightarrow\quad
 \ell_{\min}(g)^2\leq\pi\operatorname{Area}(S^2,g),
\]
with equality if and only if \(g\) has constant Gaussian curvature. This global antipodally invariant assertion has a narrower equality class than the unrestricted near-round pinched theorem. Under area \(4\pi\), its constant-curvature equality has curvature one.

The source appends higher-genus background. It attributes to Gromov the extracted bound
\[
 \operatorname{sys}(M,g)\leq C\log(\operatorname{genus}(M))
          \sqrt{\operatorname{Area}(M,g)},
\]
claiming sharp genus dependence, and to Balacheff--Sabourau [Bal-Sab2010] a bound of the displayed form
\[
 \operatorname{dias}(M,g)\leq C\sqrt{\operatorname{genus}(M)}
          \sqrt{\operatorname{Area}(M,g)},
\]
also claiming optimal genus dependence. These are retained as source-reported background, not certified sharp formulations: the first displayed genus normalization may be incomplete, the second has a damaged parenthesis, and the definition of diastole ``by (a)'' is missing. Higher-genus systole and diastole must not be silently identified with \(\ell_{\min}\) on the sphere. These auxiliary qualifications do not change the precise local sphere question above.
\subsection{Short English statement}
The aid reports a sharp pinched-curvature theorem that answers Babenko's question for metrics sufficiently close to round in \(C^2\): area \(4\pi\) gives a nonconstant closed geodesic of length at most \(2\pi\). A separate global antipodal-symmetry case has equality only at constant curvature.
\subsection{Sources}
\begin{itemize}
\item Permanent UnsolvedMath ID \(20001919\), exact public code catalogue entry; local registry \texttt{aim-geometry-notes.json}, \texttt{notes[256]}, canonical tag question.
\item AIM workshop ``Geodesics,'' \emph{Open Problems and Questions about Geodesics}, extracted item 8.7, Question 8.7.2, attributed to Babenko: \url{https://aimath.org/pastworkshops/geodesicsproblems.pdf}. The supplied verification note identifies the sphere systole as least nonconstant closed-geodesic length. No PDF was accessed.
\item The aid attributes the pinched-sphere theorem to Abbondandolo, Bramham, Hryniewicz and Salomão, without a full title, date or theorem citation. Original background names Gromov and Balacheff--Sabourau [Bal-Sab2010]; no complete bibliography or missing diastole definition is supplied.
\item Research attempt 1, ``Sharp local systolic inequality near the round two-sphere,'' labelled full\_solution, solved, proved and low novelty confidence, with classification SOLVED-IN-LITERATURE. These labels refer to the supplied \(C^2\) interpretation and are not an independently checked publication assessment.
\end{itemize}
\end{document}
